\documentclass[11pt,a4paper,oneside,openany]{book}
\usepackage{biblioflow-ista}

\renewcommand{\titolodocumento}{Impact Analysis}
\renewcommand{\sottotitolodocumento}{Impact Analysis}

% Insieme di file in forma estesa: \insieme{SIS}{riga \\ riga ...}
\newcommand{\insieme}[2]{%
  \[ \mathit{#1} = \left\{ \begin{array}{l} #2 \end{array} \right\} \]}
\newcommand{\elem}[1]{\text{\file{#1}}}

\begin{document}
\frontmatter
\copertina
\tableofcontents
\listoffigures
\listoftables

\mainmatter
\chapter{Introduzione}

BiblioFlow è una piattaforma web per la prenotazione dei posti studio e la
gestione dei servizi di una biblioteca universitaria. L'architettura combina
un'interfaccia Next.js installabile come PWA, un insieme di route handler che
accedono a PostgreSQL tramite Prisma e un servizio di automazioni periodiche
che gestisce promemoria e posti non occupati.

Questo documento di Impact Analysis si propone di esaminare in via preventiva
le modifiche richieste (Change Requests) per l'evoluzione della piattaforma.
Lo scopo dell'analisi è stimare l'estensione degli interventi necessari sui
moduli esistenti, permettendo così di pianificare lo sviluppo in modo accurato
e di mitigare i rischi legati a potenziali effetti collaterali o regressioni
nel codice. L'indagine viene svolta attraverso un'ispezione statica dei
sorgenti della versione di partenza, volta a tracciare le dipendenze logiche
tra i componenti.

\section{Obiettivi del documento}

L'analisi persegue i seguenti obiettivi operativi:
\begin{itemize}
  \item \textbf{Mappatura Architetturale:} ricostruire la struttura delle
        dipendenze di BiblioFlow attraverso l'analisi statica, identificando i
        moduli critici e le interconnessioni tra pagine, route handler e
        servizi.
  \item \textbf{Stima dell'Impatto:} determinare per ogni Change Request
        l'insieme dei file che necessitano di modifica diretta e quelli esposti
        a rischio di propagazione.
  \item \textbf{Giustificazione Formale:} fornire evidenze che motivino
        l'inclusione dei vari componenti negli insiemi di impatto previsti.
  \item \textbf{Verifica a Posteriori:} definire i criteri quantitativi
        (metriche di Precision e Recall) per valutare, a fine implementazione,
        la discrepanza tra l'impatto ipotizzato in questa sede e le modifiche
        effettivamente realizzate.
\end{itemize}

\section{Change Request Analizzate}

Le richieste di modifica analizzate sono:

{\small
\begin{xltabular}{\linewidth}{L{1cm}L{3.4cm}YL{1.3cm}}
  \toprule
  \intestazione \tb{ID} & \tb{Titolo} & \tb{Descrizione Sintetica} & \tb{Priorità} \\
  \midrule
  CR01 & Prenotazione sicura e atomica dei posti studio &
  Identità e titolarità derivate dalla sessione autenticata, regole sugli
  orari centralizzate e creazione atomica delle prenotazioni, con un vincolo
  del database che impedisce due prenotazioni attive sovrapposte sullo stesso
  posto. & Alta \\
  \midrule
  CR02 & Lista d'attesa con riassegnazione automatica dei posti &
  Coda per posto e fascia oraria quando il posto è occupato; promozione
  automatica del primo in coda quando il posto si libera per cancellazione,
  no-show o promozione non confermata, con notifiche e tracciabilità. & Media \\
  \bottomrule
  \caption{Tabella riassuntiva delle Change Request}\label{tab:riassunto}\\
\end{xltabular}}

\section{Terminologia e Acronimi}

Per classificare i componenti software in base al loro coinvolgimento nelle
modifiche, si adottano le seguenti definizioni standard:
\begin{itemize}
  \item \textbf{SIS (Starting Impact Set):} identifica i componenti che saranno
        oggetto di modifica diretta per soddisfare i requisiti della Change
        Request.
  \item \textbf{CIS (Candidate Impact Set):} comprende l'insieme esteso dei
        componenti che, pur non essendo modificati direttamente, potrebbero
        subire effetti indesiderati a causa di relazioni di dipendenza (import,
        chiamate HTTP, dati condivisi) con gli elementi del SIS.
  \item \textbf{AIS (Actual Impact Set):} rappresenta l'insieme dei componenti
        effettivamente alterati al termine della fase di implementazione.
  \item \textbf{FPIS (False Positive Impact Set):} sottoinsieme del CIS che,
        contrariamente alle previsioni, non ha richiesto interventi.
  \item \textbf{DIS (Discovered Impact Set):} componenti che hanno richiesto
        modifiche impreviste durante lo sviluppo, sfuggiti all'analisi
        iniziale.
\end{itemize}

\chapter{Change Request 01}

\section{Descrizione della Change Request}

{\small
\begin{xltabular}{\linewidth}{SY}
  \toprule
  Titolo & Prenotazione sicura e atomica dei posti studio \\
  \midrule
  Situazione Attuale &
  La creazione di una prenotazione si fida dell'identificativo utente inviato
  dal client; lettura, cancellazione ed estensione non verificano la
  titolarità. Il controllo di disponibilità e la scrittura sono operazioni
  separate e il database non impedisce sovrapposizioni. Le regole sugli orari
  sono duplicate tra creazione ed estensione e accettano intervalli invertiti;
  check-in ed estensione confrontano date e orari in rappresentazioni
  incompatibili. \\
  \midrule
  Situazione Desiderata &
  Identità e titolarità derivate dalla sessione, con risposte 401, 403 e 404
  coerenti anche sulle route dello staff. Un servizio di dominio unico per la
  validazione degli intervalli. Creazione atomica in transazione, garantita da
  un vincolo di esclusione del database; il conflitto produce un 409
  comprensibile. \\
  \bottomrule
  \caption{Dettaglio Change Request CR01}\label{tab:cr01}\\
\end{xltabular}}

\section{Analisi Statica e Architetturale}

L'analisi statica del sistema attuale evidenzia una catena di esecuzione corta
e lineare. Le pagine dello studente chiamano direttamente i route handler con
\texttt{fetch}; prima di ogni route il middleware verifica soltanto che esista
un cookie di sessione, senza leggerne il contenuto. Ciascun route handler
contiene per intero la propria logica: legge e scrive con Prisma, e solo alcune
route (check-in e area staff) interrogano \file{lib/auth.ts} per conoscere
l'utente. Non esiste uno strato di servizio condiviso tra le route di
prenotazione: ognuna replica i propri controlli.

\subsection{Call Graph Architetturale (Stato Attuale)}

Di seguito è riportata la struttura delle chiamate tra i moduli coinvolti nel
flusso di prenotazione, ricostruita tramite analisi del codice sorgente della
versione di partenza: gli archi rappresentano un import, una chiamata HTTP di
una pagina verso una route o il passaggio obbligato dal middleware.

\diagramma[0.95\linewidth]{ia-cr01-as-is}{Architettura Attuale (As-Is) del flusso di prenotazione.}

\subsection{Matrice di Raggiungibilità}

La matrice seguente rappresenta la raggiungibilità tra i moduli del grafo
precedente. I colori rappresentano la distanza topologica, cioè il numero di
archi necessari per raggiungere un componente; il trattino indica che la
destinazione non è raggiungibile.

% Generato da diagrammi/grafi_ia.py: non modificare a mano.
\begin{table}[H]
\centering
\scriptsize
\setlength{\tabcolsep}{3pt}
\begin{tabular}{lccccccccccccccc}
\toprule
\tb{Sorgente $\backslash$ Destinazione} & \rotatebox{90}{\file{prenota/page}} & \rotatebox{90}{\file{prenotazioni/page}} & \rotatebox{90}{\file{estendi/page}} & \rotatebox{90}{\file{qrcode-checkin}} & \rotatebox{90}{\file{middleware}} & \rotatebox{90}{\file{api/prenotazioni}} & \rotatebox{90}{\file{api/pren./[id]}} & \rotatebox{90}{\file{api/.../check-in}} & \rotatebox{90}{\file{api/.../estendi}} & \rotatebox{90}{\file{api/admin/posti/[id]}} & \rotatebox{90}{\file{api/admin/richieste}} & \rotatebox{90}{\file{api/admin/scanner}} & \rotatebox{90}{\file{lib/auth}} & \rotatebox{90}{\file{lib/prisma}} & \rotatebox{90}{\file{schema.prisma}} \\
\midrule
\file{prenota/page} & \cellcolor{viola!55}0 & -- & -- & -- & -- & \cellcolor{viola!38}1 & -- & -- & -- & -- & -- & -- & -- & \cellcolor{viola!26}2 & \cellcolor{viola!16}3 \\
\file{prenotazioni/page} & -- & \cellcolor{viola!55}0 & -- & -- & -- & \cellcolor{viola!38}1 & \cellcolor{viola!38}1 & -- & -- & -- & -- & -- & -- & \cellcolor{viola!26}2 & \cellcolor{viola!16}3 \\
\file{estendi/page} & -- & -- & \cellcolor{viola!55}0 & -- & -- & -- & -- & -- & \cellcolor{viola!38}1 & -- & -- & -- & -- & \cellcolor{viola!26}2 & \cellcolor{viola!16}3 \\
\file{qrcode-checkin} & -- & -- & -- & \cellcolor{viola!55}0 & -- & -- & -- & \cellcolor{viola!38}1 & -- & -- & -- & -- & \cellcolor{viola!26}2 & \cellcolor{viola!26}2 & \cellcolor{viola!16}3 \\
\file{middleware} & -- & -- & -- & -- & \cellcolor{viola!55}0 & \cellcolor{viola!38}1 & \cellcolor{viola!38}1 & \cellcolor{viola!38}1 & \cellcolor{viola!38}1 & \cellcolor{viola!38}1 & \cellcolor{viola!38}1 & \cellcolor{viola!38}1 & \cellcolor{viola!26}2 & \cellcolor{viola!26}2 & \cellcolor{viola!16}3 \\
\file{api/prenotazioni} & -- & -- & -- & -- & -- & \cellcolor{viola!55}0 & -- & -- & -- & -- & -- & -- & -- & \cellcolor{viola!38}1 & \cellcolor{viola!26}2 \\
\file{api/pren./[id]} & -- & -- & -- & -- & -- & -- & \cellcolor{viola!55}0 & -- & -- & -- & -- & -- & -- & \cellcolor{viola!38}1 & \cellcolor{viola!26}2 \\
\file{api/.../check-in} & -- & -- & -- & -- & -- & -- & -- & \cellcolor{viola!55}0 & -- & -- & -- & -- & \cellcolor{viola!38}1 & \cellcolor{viola!38}1 & \cellcolor{viola!26}2 \\
\file{api/.../estendi} & -- & -- & -- & -- & -- & -- & -- & -- & \cellcolor{viola!55}0 & -- & -- & -- & -- & \cellcolor{viola!38}1 & \cellcolor{viola!26}2 \\
\file{api/admin/posti/[id]} & -- & -- & -- & -- & -- & -- & -- & -- & -- & \cellcolor{viola!55}0 & -- & -- & \cellcolor{viola!38}1 & \cellcolor{viola!38}1 & \cellcolor{viola!26}2 \\
\file{api/admin/richieste} & -- & -- & -- & -- & -- & -- & -- & -- & -- & -- & \cellcolor{viola!55}0 & -- & -- & \cellcolor{viola!38}1 & \cellcolor{viola!26}2 \\
\file{api/admin/scanner} & -- & -- & -- & -- & -- & -- & -- & -- & -- & -- & -- & \cellcolor{viola!55}0 & \cellcolor{viola!38}1 & \cellcolor{viola!38}1 & \cellcolor{viola!26}2 \\
\file{lib/auth} & -- & -- & -- & -- & -- & -- & -- & -- & -- & -- & -- & -- & \cellcolor{viola!55}0 & \cellcolor{viola!38}1 & \cellcolor{viola!26}2 \\
\file{lib/prisma} & -- & -- & -- & -- & -- & -- & -- & -- & -- & -- & -- & -- & -- & \cellcolor{viola!55}0 & \cellcolor{viola!38}1 \\
\file{schema.prisma} & -- & -- & -- & -- & -- & -- & -- & -- & -- & -- & -- & -- & -- & -- & \cellcolor{viola!55}0 \\
\bottomrule
\end{tabular}
\caption{Matrice di Raggiungibilità tra le componenti del sistema (CR01)}\label{tab:raggiungibilita-cr01}
\end{table}


La matrice mette in evidenza due punti di concentrazione: \file{lib/prisma.ts}
e lo schema sono raggiungibili da ogni route a distanza 1 e 2, quindi una
modifica allo schema si propaga a tutto il sottosistema; il middleware
raggiunge tutte le route a distanza 1, quindi un cambiamento nel modo di
riconoscere l'utente lo coinvolge necessariamente.

\section{Impact Analysis}

\subsection{Start Impact Set (SIS)}

Lo Start Impact Set (SIS) identifica le entità software che richiedono una
modifica diretta per implementare i requisiti descritti nella Change Request.
Per la CR01 sono stati individuati sette file che costituiscono il punto di
iniezione della modifica.

{\small
\begin{xltabular}{\linewidth}{L{4.6cm}L{2cm}Y}
  \toprule
  \intestazione \tb{File (Percorso)} & \tb{Ruolo} & \tb{Motivazione dell'Inclusione} \\
  \midrule
  \endhead
  \file{prisma/schema.prisma} & Modello dati &
  \textbf{Garanzia sul database:} la sovrapposizione è oggi controllata solo
  dal codice. Serve un vincolo che renda impossibili due prenotazioni attive
  sovrapposte sullo stesso posto. \\
  \file{src/lib/auth.ts} & Sicurezza &
  \textbf{Identità dalla sessione:} mancano funzioni comuni per ottenere
  l'utente autenticato e verificarne ruolo e titolarità; ogni route le
  reimplementa o le omette. \\
  \file{api/prenotazioni/route.ts} & Controller &
  \textbf{Creazione atomica:} legge \texttt{userId} dal corpo della richiesta
  ed esegue verifica e scrittura in due momenti distinti; l'elenco accetta un
  \texttt{userId} arbitrario come filtro. \\
  \file{api/prenotazioni/[id]/route.ts} & Controller &
  \textbf{Titolarità:} lettura, cancellazione e check-in tramite
  \texttt{PATCH} non verificano che la prenotazione appartenga all'utente. \\
  \file{api/prenotazioni/[id]/check-in/route.ts} & Controller &
  \textbf{Confronto data/ora:} confronta l'istante della richiesta con un
  orario \texttt{@db.Time} ancorato al 1970 e rifiuta check-in validi. \\
  \file{api/prenotazioni/[id]/estendi/route.ts} & Controller &
  \textbf{Validazione duplicata:} non verifica la sessione, replica le regole
  sugli orari e calcola la durata tra date incompatibili. \\
  \file{app/prenota/page.tsx} & Interfaccia &
  \textbf{Contratto della richiesta:} invia l'identificativo dell'utente nel
  corpo della prenotazione; il contratto con l'API deve cambiare. \\
  \bottomrule
  \caption{Start Impact Set (SIS) per la CR01}\label{tab:sis-cr01}\\
\end{xltabular}}

\insieme{SIS}{
  \elem{prisma/schema.prisma},\ \elem{src/lib/auth.ts},\\
  \elem{api/prenotazioni/route.ts},\ \elem{api/prenotazioni/[id]/route.ts},\\
  \elem{api/prenotazioni/[id]/check-in/route.ts},\\
  \elem{api/prenotazioni/[id]/estendi/route.ts},\ \elem{app/prenota/page.tsx}}

\subsection{Candidate Impact Set (CIS)}

L'analisi del Candidate Impact Set (CIS) mira a identificare i componenti che
dipendono, direttamente o indirettamente, dagli elementi del SIS.
Inizialmente l'indagine è stata condotta tramite ricerca testuale degli import
e delle chiamate \texttt{fetch} nel codice sorgente, con il comando
\texttt{grep}. Data la natura puramente sintattica di questo metodo, si è
proceduto a una revisione manuale delle dipendenze, appoggiata al diagramma dei
componenti e alla matrice ruoli-operazioni redatti sulla versione di partenza
prima di intervenire sul codice.

A differenza di quanto avviene per un intervento isolato, il CIS non coincide
con il SIS: la revisione ha aggiunto sei componenti.

{\small
\begin{xltabular}{\linewidth}{L{4.6cm}Y}
  \toprule
  \intestazione \tb{File (Percorso)} & \tb{Motivo della Propagazione} \\
  \midrule
  \endhead
  \file{src/middleware.ts} &
  Precede tutte le route e oggi riconosce l'utente dalla sola presenza del
  cookie: una gestione dell'identità basata sulla sessione potrebbe
  richiederne l'adeguamento. \\
  \file{api/admin/posti/[id]/route.ts} &
  Restituisce le prenotazioni di un posto; la matrice ruoli-operazioni rileva
  che verifica la sessione ma non il ruolo dello staff. \\
  \file{api/admin/richieste/route.ts} &
  Route dello staff che non verifica né la sessione né il ruolo: rientra nel
  requisito di risposte coerenti per l'area staff. \\
  \file{api/admin/scanner/validate/route.ts} &
  Esegue il check-in tramite QR e scrive gli stessi stati della prenotazione
  toccati dalla CR. \\
  \file{app/prenotazioni/page.tsx} &
  Consuma l'elenco e le operazioni sulla singola prenotazione, il cui
  contratto (filtri, codici di errore) può cambiare. \\
  \file{app/prenotazioni/[id]/estendi/page.tsx} &
  Consuma la route di estensione, le cui regole e risposte saranno
  centralizzate. \\
  \bottomrule
  \caption{Componenti aggiunti dal Candidate Impact Set per la CR01}\label{tab:cis-cr01}\\
\end{xltabular}}

\insieme{CIS}{
  \mathit{SIS}\ \cup\ \{\elem{src/middleware.ts},\ \elem{api/admin/posti/[id]/route.ts},\\
  \quad \elem{api/admin/richieste/route.ts},\ \elem{api/admin/scanner/validate/route.ts},\\
  \quad \elem{app/prenotazioni/page.tsx},\ \elem{app/prenotazioni/[id]/estendi/page.tsx}\}}

\section{Implementazione}

\subsection{Actual Impact Set (AIS)}

Durante la fase di implementazione, la scelta di raccogliere le regole di
prenotazione in un servizio di dominio, invece di correggerle route per route,
ha portato a una divergenza tra l'Impact Set stimato (CIS) e quello effettivo
(AIS). L'AIS comprende i seguenti elementi:
\begin{itemize}
  \item \textbf{File Modificati (dal CIS):}
  \begin{itemize}
    \item \file{prisma/schema.prisma}: aggiornato per i nuovi indici sulle
          prenotazioni attive.
    \item \file{src/lib/auth.ts}: aggiunte \texttt{requireUser},
          \texttt{requireRole} e \texttt{assertOwnership}, con un errore di
          autorizzazione tipizzato che le route traducono in 401, 403 o 404.
    \item \file{api/prenotazioni/route.ts},
          \file{api/prenotazioni/[id]/route.ts},
          \file{api/prenotazioni/[id]/check-in/route.ts},
          \file{api/prenotazioni/[id]/estendi/route.ts}: identità dalla
          sessione, titolarità, delega al servizio di dominio e traduzione
          degli errori in risposte HTTP.
    \item \file{api/admin/posti/[id]/route.ts},
          \file{api/admin/richieste/route.ts}: verifica del ruolo dello staff e
          validazione dei dati in ingresso.
  \end{itemize}
  \item \textbf{Nuovi File Creati (Estensione):}
  \begin{itemize}
    \item \file{src/lib/prenotazioni-service.ts}: servizio di dominio con la
          validazione degli intervalli e la creazione in transazione
          \texttt{Serializable}.
    \item \file{src/lib/prenotazioni-errors.ts}: errori di dominio con codice
          stabile e messaggio comprensibile, tra cui il conflitto di
          disponibilità.
    \item \file{prisma/migrations/..._vincolo_anti_sovrapposizione_prenotazioni}:
          vincolo di esclusione PostgreSQL sugli intervalli delle prenotazioni
          attive dello stesso posto.
  \end{itemize}
\end{itemize}

Formalmente, l'insieme dei file impattati risulta essere:

\insieme{AIS_{CR01}}{
  \elem{prisma/schema.prisma},\ \elem{src/lib/auth.ts},\\
  \elem{api/prenotazioni/route.ts},\ \elem{api/prenotazioni/[id]/route.ts},\\
  \elem{api/prenotazioni/[id]/check-in/route.ts},\\
  \elem{api/prenotazioni/[id]/estendi/route.ts},\\
  \elem{api/admin/posti/[id]/route.ts},\ \elem{api/admin/richieste/route.ts},\\
  \elem{src/lib/prenotazioni-service.ts}\ \text{(New)},\\
  \elem{src/lib/prenotazioni-errors.ts}\ \text{(New)},\\
  \text{migrazione del vincolo anti-sovrapposizione (New)}}

La differenza tra stima ed effettivo evidenzia un approccio additivo per la
logica di dominio: le regole sono state spostate in moduli nuovi e le route si
sono ridotte a tradurre input ed errori, invece di essere corrette una per
una.

\subsection{False Positive Impact Set (FPIS)}

Il False Positive Impact Set (FPIS) è costituito dai componenti che erano
stati stimati come impattati dalla modifica (membri del CIS), ma che non sono
stati modificati durante l'implementazione (non membri dell'AIS). La tabella
seguente descrive questi file e la giustificazione della loro esclusione.

{\small
\begin{xltabular}{\linewidth}{L{4.6cm}Y}
  \toprule
  \intestazione \tb{Classe/File} & \tb{Motivo dell'Esclusione} \\
  \midrule
  \endhead
  \file{app/prenota/page.tsx} &
  La pagina continua a inviare l'identificativo dell'utente, ma il server lo
  ignora e usa quello della sessione: la correzione è stata fatta interamente
  lato server, senza cambiare il contratto della pagina. \\
  \file{src/middleware.ts} &
  Le route ottengono l'utente dalla sessione per conto proprio; il
  middleware, che si limita a respingere le richieste prive di cookie, non ha
  richiesto modifiche per questa CR. \\
  \file{api/admin/scanner/validate/route.ts} &
  Verificava già sessione, ruolo e firma del QR; la CR non ha introdotto
  nuovi stati della prenotazione che lo scanner dovesse riconoscere. \\
  \file{app/prenotazioni/page.tsx} &
  I contratti di elenco e cancellazione sono rimasti compatibili: le nuove
  risposte di errore sono già gestite dal codice esistente. \\
  \file{app/prenotazioni/[id]/estendi/page.tsx} &
  La route di estensione mantiene la stessa richiesta e la stessa risposta di
  successo; cambiano solo le regole lato server. \\
  \bottomrule
  \caption{False Positive Impact Set (FPIS) per la CR01}\label{tab:fpis-cr01}\\
\end{xltabular}}

\textbf{Analisi:} i falsi positivi derivano da una stima prudente: ogni
consumatore delle route di prenotazione è stato incluso perché il contratto
poteva cambiare. Spostare la correzione lato server ha permesso di mantenere
compatibili i contratti e di lasciare intatte le pagine. Il middleware è stato
poi rivisto durante la manutenzione successiva, quando è emerso che il solo
controllo della presenza del cookie non era sufficiente: ora verifica davvero
il token di sessione.

\subsection{Discovered Impact Set (DIS)}

Il Discovered Impact Set (DIS) comprende i file che non erano stati inclusi
nell'insieme di impatto iniziale, ma che sono stati creati durante l'effettiva
implementazione della Change Request. Questi componenti rappresentano il
lavoro \enquote{scoperto} durante lo sviluppo, dovuto a decisioni
architetturali prese per migliorare la manutenibilità del codice.

Le seguenti classi compongono il DIS per la CR01:

{\small
\begin{xltabular}{\linewidth}{L{4.6cm}Y}
  \toprule
  \intestazione \tb{Classe/File} & \tb{Motivazione dell'Inclusione} \\
  \midrule
  \endhead
  \file{src/lib/prenotazioni-service.ts} &
  Invece di correggere quattro route con quattro copie delle stesse regole, è
  stato creato un servizio di dominio indipendente da richiesta e risposta
  HTTP: validazione degli intervalli, verifica della sala e del posto,
  creazione ed estensione in transazione. \\
  \file{src/lib/prenotazioni-errors.ts} &
  Le route devono distinguere i casi d'errore senza conoscere il database:
  gli errori di dominio portano un codice stabile e un messaggio
  comprensibile, che la route traduce nel codice HTTP corretto. \\
  Migrazione del vincolo anti-sovrapposizione &
  Lo schema Prisma non consente di dichiarare un vincolo di esclusione sugli
  intervalli: è stato necessario scriverlo in SQL in una migrazione dedicata,
  oltre alla modifica di \file{schema.prisma} prevista nel SIS. \\
  \bottomrule
  \caption{Discovered Impact Set (DIS) per la CR01}\label{tab:dis-cr01}\\
\end{xltabular}}

\textbf{Analisi:} la presenza di questi file nel DIS mostra che una parte
dell'impatto si è spostata dalla modifica del codice esistente alla creazione
di moduli nuovi. È lo stesso fenomeno che riduce l'FPIS: concentrando le
regole in un punto solo, i consumatori sono rimasti intatti.

\section{Metrics}

In questa sezione valutiamo l'accuratezza della stima dell'impatto utilizzando
le metriche di Precision e Recall.
\begin{itemize}
  \item \textbf{Candidate Impact Set (CIS):} 13 file
  \item \textbf{Actual Impact Set (AIS):} 11 file
  \item \textbf{Intersezione (CIS $\cap$ AIS):} 8 file
\end{itemize}

\paragraph{Precision (Precisione)}
Misura la frazione di file previsti che sono stati effettivamente modificati.
\[ \mathit{Precision} = \frac{|\mathit{CIS} \cap \mathit{AIS}|}{|\mathit{CIS}|} = \frac{8}{13} \approx 0{,}62\ (61{,}5\%) \]
Un valore del 62\% indica che la maggior parte dei file previsti è stata
coinvolta; i cinque falsi positivi sono consumatori il cui contratto è
rimasto compatibile.

\paragraph{Recall (Richiamo)}
Misura la frazione di file effettivamente modificati che erano stati previsti
correttamente all'inizio.
\[ \mathit{Recall} = \frac{|\mathit{CIS} \cap \mathit{AIS}|}{|\mathit{AIS}|} = \frac{8}{11} \approx 0{,}73\ (72{,}7\%) \]
Un valore del 73\% indica che quasi tre quarti del lavoro effettivo erano
stati previsti; la parte restante corrisponde ai tre file nuovi del DIS, nati
dalla scelta di introdurre un servizio di dominio.

\chapter{Change Request 02}

\section{Descrizione della Change Request}

{\small
\begin{xltabular}{\linewidth}{SY}
  \toprule
  Titolo & Lista d'attesa con riassegnazione automatica dei posti \\
  \midrule
  Situazione Attuale &
  Se il posto è occupato nella fascia scelta lo studente riceve solo un
  errore. Quando una prenotazione viene annullata o il titolare non si
  presenta, il posto torna disponibile senza avvisare nessuno; l'automazione
  periodica si limita a segnare il no-show e a liberare il posto. \\
  \midrule
  Situazione Desiderata &
  Ingresso in coda per posto e fascia dalla mappa, dalla griglia mobile o dal
  messaggio di conflitto; ordine di arrivo senza duplicati; promozione
  automatica del primo in coda su cancellazione (studente o staff), no-show e
  promozione non confermata, con notifica, scadenza della promozione e
  tracciabilità completa. \\
  \bottomrule
  \caption{Dettaglio Change Request CR02}\label{tab:cr02}\\
\end{xltabular}}

\section{Analisi Statica e Architetturale (Liberazione dei Posti)}

L'architettura coinvolta dalla CR02 è più distribuita di quella della CR01:
un posto può liberarsi per tre vie indipendenti (cancellazione dello studente,
cancellazione dello staff, automazione periodica) che oggi non condividono
alcun codice. L'analisi statica ha inoltre ricostruito la catena dei
componenti di aggiornamento in tempo reale, presenti nel codice ma non
collegati all'interfaccia: nessuna pagina usa l'hook \file{use-sse.ts} e
nessun modulo invoca \file{realtime-events.ts}.

\subsection{Call Graph Architetturale (Stato Attuale)}

Per comprendere l'impatto della lista d'attesa è stata mappata la catena che
va dalle interfacce ai punti in cui un posto si libera, fino allo schema dei
dati. Il diagramma adotta una prospettiva orientata al flusso dei dati e
mostra il disaccoppiamento tra i tre percorsi di liberazione e il canale in
tempo reale.

\diagramma[0.82\linewidth]{ia-cr02-as-is}{Liberazione dei posti e canale realtime As-Is.}

\subsection{Matrice di Raggiungibilità (Cross-Layer)}

La seguente matrice quantifica la distanza logica tra i livelli
dell'architettura. Nel contesto della CR02 è significativo ciò che
\emph{non} è raggiungibile: dall'interfaccia di prenotazione non si arriva al
servizio delle automazioni né al canale SSE, quindi l'esito di una promozione
non ha oggi alcuna strada per tornare all'utente se non tramite la base di dati
e le notifiche.

% Generato da diagrammi/grafi_ia.py: non modificare a mano.
\begin{table}[H]
\centering
\footnotesize
\setlength{\tabcolsep}{3pt}
\begin{tabular}{lcccccccc}
\toprule
\tb{Sorgente $\backslash$ Destinazione} & \rotatebox{90}{UI prenotazione} & \rotatebox{90}{UI staff} & \rotatebox{90}{Client realtime} & \rotatebox{90}{Route handler} & \rotatebox{90}{Cron} & \rotatebox{90}{automation-service} & \rotatebox{90}{Canale SSE} & \rotatebox{90}{Schema dati} \\
\midrule
UI prenotazione & \cellcolor{viola!55}0 & -- & -- & \cellcolor{viola!38}1 & -- & -- & -- & \cellcolor{viola!26}2 \\
UI staff & -- & \cellcolor{viola!55}0 & -- & \cellcolor{viola!38}1 & -- & -- & -- & \cellcolor{viola!38}1 \\
Client realtime & -- & -- & \cellcolor{viola!55}0 & -- & -- & -- & \cellcolor{viola!38}1 & -- \\
Route handler & -- & -- & -- & \cellcolor{viola!55}0 & -- & -- & -- & \cellcolor{viola!38}1 \\
Cron & -- & -- & -- & -- & \cellcolor{viola!55}0 & \cellcolor{viola!38}1 & -- & \cellcolor{viola!26}2 \\
automation-service & -- & -- & -- & -- & -- & \cellcolor{viola!55}0 & -- & \cellcolor{viola!38}1 \\
Canale SSE & -- & -- & -- & -- & -- & -- & \cellcolor{viola!55}0 & \cellcolor{viola!38}1 \\
Schema dati & -- & -- & -- & -- & -- & -- & -- & \cellcolor{viola!55}0 \\
\bottomrule
\end{tabular}
\caption{Matrice di Raggiungibilità Cross-Layer (CR02)}\label{tab:raggiungibilita-cr02}
\end{table}


\section{Impact Analysis}

\subsection{Start Impact Set (SIS)}

Lo Start Impact Set per la CR02 comprende i componenti che possiedono i punti
in cui un posto si libera e il modello dei dati che deve rappresentare la
coda.

{\small
\begin{xltabular}{\linewidth}{L{4.6cm}L{2cm}Y}
  \toprule
  \intestazione \tb{File (Percorso)} & \tb{Tipo} & \tb{Motivazione Specifica} \\
  \midrule
  \endhead
  \file{prisma/schema.prisma} & Modello dati &
  \textbf{Nuova entità:} la richiesta in coda deve essere persistita con
  utente, posto, intervallo, stato e ordine di arrivo, impedendo i duplicati
  attivi. \\
  \file{src/lib/automation-service.ts} & Servizio &
  \textbf{No-show e scadenze:} il rilascio automatico dei posti non occupati è
  uno dei tre eventi di liberazione; qui va avviata la promozione e gestita la
  scadenza delle promozioni non confermate. \\
  \file{api/admin/prenotazioni/route.ts} & Controller &
  \textbf{Cancellazione dello staff:} l'annullamento di una prenotazione da
  parte dello staff libera il posto e deve tentare la promozione. \\
  \file{api/prenotazioni/[id]/route.ts} & Controller &
  \textbf{Cancellazione dello studente:} l'annullamento da parte del titolare
  è il caso più frequente di posto liberato. \\
  \bottomrule
  \caption{Start Impact Set (SIS) Dettagliato per CR02}\label{tab:sis-cr02}\\
\end{xltabular}}

\insieme{SIS}{
  \elem{prisma/schema.prisma},\ \elem{src/lib/automation-service.ts},\\
  \elem{api/admin/prenotazioni/route.ts},\ \elem{api/prenotazioni/[id]/route.ts}}

\subsection{Candidate Impact Set (CIS)}

Il Candidate Impact Set identifica i componenti che potrebbero subire effetti
collaterali a causa delle modifiche introdotte nel SIS. Come per la CR01,
l'indagine è partita dalla ricerca testuale di import, chiamate e tipi
condivisi (stati della prenotazione, tipi di notifica e di evento) ed è stata
completata con una revisione manuale. A differenza di un intervento isolato,
la propagazione è ampia: i nuovi stati ed eventi attraversano statistiche,
anomalie, interfacce e automazioni.

{\small
\begin{xltabular}{\linewidth}{L{4.6cm}Y}
  \toprule
  \intestazione \tb{File (Percorso)} & \tb{Motivo della Propagazione} \\
  \midrule
  \endhead
  \file{app/prenota/page.tsx},
  \file{components/mappa-biblioteca.tsx},
  \file{components/mobile-posti-grid.tsx} &
  Sono i punti in cui lo studente incontra il posto occupato: la proposta di
  entrare in coda deve comparire qui. \\
  \file{api/cron/automations/route.ts} &
  Invoca il servizio delle automazioni: esecuzioni concorrenti potrebbero
  promuovere due volte lo stesso utente. \\
  \file{api/admin/anomalie/route.ts},
  \file{app/admin/page.tsx} &
  Contano no-show e cancellazioni per dashboard e anomalie: i nuovi stati ed
  eventi potrebbero alterarne i valori. \\
  \file{api/admin/statistiche/route.ts},
  \file{app/admin/statistiche/page.tsx} &
  Le statistiche aggregano gli stati delle prenotazioni; la coda richiede un
  indicatore proprio senza inquinare quelli esistenti. \\
  \file{api/admin/scanner/validate/route.ts} &
  Esegue il check-in: deve accettare le prenotazioni nate da una promozione e
  non le richieste ancora in coda. \\
  \file{api/disponibilita-giorni/route.ts} &
  Calcola i giorni prenotabili dagli stati delle prenotazioni. \\
  \file{lib/realtime-events.ts}, \file{lib/sse-emitter.ts},
  \file{api/sse/posti/route.ts}, \file{hooks/use-sse.ts} &
  Canale previsto per avvisare in tempo reale dei posti liberati e delle
  promozioni. \\
  \bottomrule
  \caption{Componenti aggiunti dal Candidate Impact Set per la CR02}\label{tab:cis-cr02}\\
\end{xltabular}}

Il CIS conta quindi 18 file: i 4 del SIS e i 14 elencati.

\section{Implementazione}

\subsection{Actual Impact Set (AIS)}

Durante l'implementazione della CR02 le scelte architetturali hanno
privilegiato il riuso del servizio di dominio creato per la CR01: ingresso in
coda, annullamento, posizione e promozione sono funzioni dello stesso
servizio, chiamate da tutti i punti di liberazione. L'AIS comprende i
seguenti elementi:
\begin{itemize}
  \item \textbf{File Modificati (dal CIS):}
  \begin{itemize}
    \item \file{prisma/schema.prisma}: entità \texttt{ListaAttesa} con stati
          \texttt{IN\_ATTESA}, \texttt{PROMOSSA}, \texttt{SCADUTA},
          \texttt{ANNULLATA} e nuovi tipi di notifica e di evento.
    \item \file{src/lib/automation-service.ts}: promozione dopo il no-show,
          scadenza delle promozioni non confermate entro quindici minuti,
          protezione dalle doppie esecuzioni.
    \item \file{api/cron/automations/route.ts}: lock di PostgreSQL che
          impedisce due esecuzioni contemporanee.
    \item \file{api/admin/prenotazioni/route.ts}: promozione dopo
          l'annullamento dello staff ed esito restituito all'interfaccia.
    \item \file{app/prenota/page.tsx}, \file{mappa-biblioteca.tsx},
          \file{mobile-posti-grid.tsx}: proposta di ingresso in coda sul posto
          occupato e nel messaggio di conflitto.
    \item \file{api/admin/statistiche/route.ts},
          \file{app/admin/statistiche/page.tsx}: indicatore dedicato alla
          coda.
    \item \file{lib/realtime-events.ts}, \file{api/sse/posti/route.ts},
          \file{hooks/use-sse.ts}: nuovi eventi di coda e promozione sul canale
          in tempo reale.
  \end{itemize}
  \item \textbf{Nuovi File Creati o Modificati fuori stima:}
  \begin{itemize}
    \item \file{src/lib/prenotazioni-service.ts}: nuove funzioni di
          dominio per la coda (\texttt{entraInCoda},
          \texttt{annullaRichiestaCoda}, \texttt{posizioneInCoda},
          \texttt{promuoviPrimoInCoda}).
    \item \file{api/prenotazioni/coda/route.ts} (New): elenco, ingresso e
          uscita dalla coda per lo studente.
    \item \file{hooks/use-ingresso-coda.ts} (New) e
          \file{lib/eventi-coda.ts} (New): logica client condivisa tra mappa e
          griglia, tipi degli eventi di coda.
    \item \file{app/notifiche/tipo-config.tsx} (New) e
          \file{app/notifiche/page.tsx}: resa dei nuovi tipi di notifica.
    \item \file{src/middleware.ts}, \file{admin/prenotazioni-actions.tsx},
          \file{admin/statistiche-charts.tsx}: protezione della route della
          coda, esito della promozione per lo staff, grafico della coda.
    \item Due migrazioni (New): tabella della lista d'attesa con indici e
          vincolo di unicità delle richieste attive; estensione degli enum.
  \end{itemize}
\end{itemize}

Formalmente, l'insieme dei file impattati risulta essere:

\insieme{AIS_{CR02}}{
  \elem{prisma/schema.prisma},\ \elem{src/lib/automation-service.ts},\\
  \elem{api/cron/automations/route.ts},\ \elem{api/admin/prenotazioni/route.ts},\\
  \elem{app/prenota/page.tsx},\ \elem{mappa-biblioteca.tsx},\ \elem{mobile-posti-grid.tsx},\\
  \elem{api/admin/statistiche/route.ts},\ \elem{app/admin/statistiche/page.tsx},\\
  \elem{lib/realtime-events.ts},\ \elem{api/sse/posti/route.ts},\ \elem{hooks/use-sse.ts},\\
  \elem{src/lib/prenotazioni-service.ts},\ \elem{app/notifiche/page.tsx},\\
  \elem{src/middleware.ts},\ \elem{admin/prenotazioni-actions.tsx},\\
  \elem{admin/statistiche-charts.tsx},\ \elem{api/prenotazioni/coda/route.ts}\ \text{(New)},\\
  \elem{hooks/use-ingresso-coda.ts}\ \text{(New)},\ \elem{lib/eventi-coda.ts}\ \text{(New)},\\
  \elem{app/notifiche/tipo-config.tsx}\ \text{(New)},\\
  \text{migrazione della lista d'attesa (New)},\\
  \text{migrazione di estensione degli enum (New)}}

\subsection{False Positive Impact Set (FPIS)}

Il False Positive Impact Set evidenzia i componenti inclusi nel CIS che sono
rimasti inalterati durante l'implementazione.

{\small
\begin{xltabular}{\linewidth}{L{4.6cm}Y}
  \toprule
  \intestazione \tb{Classe/File} & \tb{Motivo dell'Esclusione} \\
  \midrule
  \endhead
  \file{api/admin/anomalie/route.ts}, \file{app/admin/page.tsx} &
  La coda è un'entità separata: le richieste in attesa non sono prenotazioni
  e i conteggi di no-show e cancellazioni non sono cambiati. \\
  \file{api/admin/scanner/validate/route.ts} &
  Una promozione produce una normale prenotazione \texttt{CONFERMATA}: lo
  scanner la riconosce senza modifiche, e una richiesta in coda non ha un QR. \\
  \file{lib/sse-emitter.ts} &
  I nuovi eventi sono stati aggiunti in \file{realtime-events.ts}; il
  meccanismo di consegna è rimasto invariato. \\
  \file{api/disponibilita-giorni/route.ts} &
  I giorni prenotabili dipendono dal calendario e dalle prenotazioni attive,
  non dalla coda. \\
  \file{api/prenotazioni/[id]/route.ts} &
  \textbf{Omissione, non vero falso positivo.} A fine CR la cancellazione da
  parte dello studente non avviava la promozione: il caso era verificato solo
  a livello di servizio. Il difetto è emerso dalla verifica dal vivo ed è stato
  corretto nella manutenzione successiva, collegando la route alla
  promozione. \\
  \bottomrule
  \caption{False Positive Impact Set (FPIS) per la CR02}\label{tab:fpis-cr02}\\
\end{xltabular}}

\textbf{Analisi:} cinque dei sei elementi sono veri falsi positivi, dovuti
alla scelta di modellare la coda come entità separata invece che come nuovo
stato della prenotazione: i componenti che leggono le prenotazioni non hanno
dovuto cambiare. Il sesto è un caso diverso e più istruttivo: la stima era
corretta, l'implementazione incompleta. Per questo l'FPIS non va letto come
misura della sola qualità della stima.

\subsection{Discovered Impact Set (DIS)}

Il Discovered Impact Set comprende i file non inclusi nel CIS che sono stati
creati o modificati durante l'implementazione.

{\small
\begin{xltabular}{\linewidth}{L{4.6cm}Y}
  \toprule
  \intestazione \tb{Classe/File} & \tb{Motivazione dell'Inclusione} \\
  \midrule
  \endhead
  \file{src/lib/prenotazioni-service.ts} &
  Le operazioni di coda devono condividere la transazione e le regole della
  prenotazione: sono state aggiunte al servizio di dominio della CR01, con
  lock sulle righe per impedire doppie promozioni concorrenti. \\
  \file{api/prenotazioni/coda/route.ts} &
  Nessuna route esistente poteva ospitare elenco, ingresso e uscita dalla
  coda senza violare la responsabilità singola: è stata creata una route
  dedicata. \\
  \file{hooks/use-ingresso-coda.ts}, \file{lib/eventi-coda.ts} &
  La stessa logica di ingresso in coda serviva a mappa, griglia mobile e
  messaggio di conflitto: è stata estratta in un hook e in un modulo di tipi
  condivisi. \\
  \file{app/notifiche/page.tsx}, \file{app/notifiche/tipo-config.tsx} &
  La pagina delle notifiche associava icone e colori solo ai tipi esistenti:
  i nuovi tipi richiedevano una configurazione con un ripiego per i tipi
  sconosciuti. \\
  \file{src/middleware.ts} &
  La nuova route della coda doveva essere protetta come le altre route
  riservate. \\
  \file{admin/prenotazioni-actions.tsx}, \file{admin/statistiche-charts.tsx} &
  Lo staff deve sapere chi è stato promosso dopo un annullamento e vedere
  l'andamento della coda. \\
  Migrazioni \texttt{lista\_attesa} ed estensione degli enum &
  PostgreSQL non permette di usare un nuovo valore di enum nella stessa
  transazione che lo aggiunge: tabella e valori sono stati separati in due
  migrazioni. \\
  \bottomrule
  \caption{Discovered Impact Set (DIS) per la CR02}\label{tab:dis-cr02}\\
\end{xltabular}}

\textbf{Analisi:} il DIS è ampio perché la CR02 ha seguito una strategia di
estensione: componenti nuovi e dedicati, invece di modificare quelli
esistenti. La verifica successiva ha mostrato un ulteriore punto sfuggito sia
alla stima sia all'implementazione: nessuna schermata permetteva allo studente
di vedere le proprie richieste in coda e di uscirne, benché l'API lo
consentisse. È stata aggiunta in seguito una scheda dedicata in \enquote{Le
mie prenotazioni}.

\subsection{Metrics}

Valutiamo ora l'accuratezza della stima dell'impatto per questa Change
Request utilizzando le metriche standard.
\begin{itemize}
  \item \textbf{Candidate Impact Set (CIS):} 18 file
  \item \textbf{Actual Impact Set (AIS):} 23 file
  \item \textbf{Intersezione (CIS $\cap$ AIS):} 12 file
\end{itemize}

\paragraph{Precision (Precisione)}
Indica quanto la stima iniziale fosse mirata, ovvero quanta parte dei file
previsti è stata effettivamente modificata.
\[ \mathit{Precision} = \frac{|\mathit{CIS} \cap \mathit{AIS}|}{|\mathit{CIS}|} = \frac{12}{18} \approx 0{,}67\ (66{,}7\%) \]
Un valore del 67\% indica una precisione discreta: la scelta di rappresentare
la coda come entità separata ha reso superflue le modifiche a statistiche,
anomalie e scanner che la stima aveva previsto per prudenza.

\paragraph{Recall (Richiamo)}
Indica quanto del lavoro finale effettivamente svolto era stato previsto
all'inizio.
\[ \mathit{Recall} = \frac{|\mathit{CIS} \cap \mathit{AIS}|}{|\mathit{AIS}|} = \frac{12}{23} \approx 0{,}52\ (52{,}2\%) \]
Un valore del 52\% conferma che circa metà dell'implementazione ha
riguardato file nuovi o componenti non previsti, come il servizio di dominio
esteso, la route della coda e l'hook condiviso. Riflette la scelta, presa in
corso d'opera, di aggiungere moduli dedicati invece di modificare quelli
esistenti.

\end{document}
